\documentclass[12pt,a4paper]{article}
\usepackage[T1]{fontenc} 
\usepackage[latin1]{inputenc}
\usepackage{graphicx}
\usepackage[usenames]{color}
\usepackage{fancybox}
\usepackage{url}
\usepackage[finnish]{babel}
\usepackage{mathrsfs,amsfonts,amsmath,amssymb}
\usepackage[noindent]{lc2005}

\addtolength{\textwidth}{2cm}
\addtolength{\hoffset}{-1cm}

\pagestyle{empty}

\begin{document}


\begin{center}
\raisebox{-0.4cm}{\includegraphics[height=1.5cm]{soihtu}} \quad 
{\Large\itshape Turun yliopiston matematiikan laitoksen}\\[1cm]

\textcolor{Periwinkle}{\rule{1\textwidth}{2pt}}\bigskip

\textcolor{Periwinkle}{\scalebox{6}{\LaTeX-}}

\textcolor{YellowOrange}{\scalebox{4}{\script kurssi}}\medskip

\textcolor{Periwinkle}{\rule{1\textwidth}{2pt}}
\end{center}

\vspace*{-5.8cm}\hspace*{12cm}
\rotatebox{45}{\shadowbox{\parbox{3.5cm}
    {10--14.1. 2005\\klo 10.15--12.00\\XXIV ja XXI}}}
\vspace*{1cm}

\begin{quote}
\large
\LaTeX{} on erityisesti matemaattisiin teksteihin
soveltuva tait\-to-\-oh\-jel\-ma ja sit� k�ytet��n esimerkiksi kirjojen,
Pro gradu -tut\-ki\-el\-mi\-en, raporttien ja esitysgrafiikan tekoon.
T�m� mainos on tehty \LaTeX illa.
\end{quote}

{\large \textbf{Kurssi on suunnattu}}
\begin{itemize}
\item opettajaopiskelijoille ja
\item oppimateriaaleja laativille,
\item aineen- ja graduntekij�ille sek�
\item muille taitto-ohjelmista kiinnostuneille
\end{itemize}
ja siit� on mahdollista saada yksi (1) yleisopintotasoinen
opintoviikko. Kaikki ne, jotka laativat matemaattisluonteista
materiaalia hy�tyv�t kurssista.\bigskip

{\large \textbf{Aikataulu ja materiaali}}

Kurssi j�rjestet��n intensiiviopetuksena 10--14.1. 2005 klo 10.15--12.00
saleissa XXIV (ma, ti) ja XXI (ke-pe).

\newcommand{\baseurl}{http://users.utu.fi/avirkki/latex2005/}
\urlstyle{sf}

\textbf{Kurssin internetsivu:} \\ 
\baseurl

\textbf{Kirja:} \\
\textsl{Helmut Kopka and Patrick W. Daly, Guide to LaTeX,
4th edition}\bigskip


\end{document}
